% TODO make hw stylesheet
\documentclass[10pt]{scrartcl}
\usepackage[a4paper, total={6in, 8in}]{geometry}
\usepackage[usenames,dvipsnames,svgnames]{xcolor}
\usepackage[shortlabels]{enumitem}
\usepackage[framemethod=TikZ]{mdframed}
\usepackage{amsmath,amssymb,amsthm}
\usepackage[colorlinks]{hyperref}
\usepackage{multicol}
\usepackage{tabularx}

\hypersetup{
	colorlinks=true,
	linkcolor=blue,
	filecolor=magenta,      
	urlcolor=magenta,
	pdftitle={MAT 362 HW}
}

\usepackage{microtype}
\usepackage{mathtools}
\usepackage[headsepline]{scrlayer-scrpage}
\usepackage{thmtools}
\usepackage[textsize=footnotesize]{todonotes}

\mdfdefinestyle{mdbluebox}{roundcorner=10pt,innerbottommargin=9pt,
	linecolor=blue,backgroundcolor=TealBlue!5,}
\declaretheoremstyle[headfont=\sffamily\bfseries\color{MidnightBlue},
mdframed={style=mdbluebox},]{thmbluebox}

\mdfdefinestyle{mdbloobox}{frametitlefont=\bfseries,innerbottommargin=8pt,
	nobreak=true,backgroundcolor=TealBlue!5,linecolor=blue,}
\declaretheoremstyle[headfont=\bfseries\color{MidnightBlue},
mdframed={style=mdbloobox},headpunct={\\[3pt]},postheadspace=0pt,]{thmbloobox}

\mdfdefinestyle{mdredbox}{frametitlefont=\bfseries,innerbottommargin=8pt,
	nobreak=true,backgroundcolor=Salmon!5,linecolor=RawSienna,}
\declaretheoremstyle[headfont=\bfseries\color{RawSienna},
mdframed={style=mdredbox},headpunct={\\[3pt]},postheadspace=0pt,]{thmredbox}

\mdfdefinestyle{mdgreenbox}{linecolor=ForestGreen,backgroundcolor=ForestGreen!5,
	linewidth=2pt,rightline=false,leftline=true,topline=false,bottomline=false,}
\declaretheoremstyle[headfont=\bfseries\sffamily\color{ForestGreen!70!black},
mdframed={style=mdgreenbox},headpunct={ --- },]{thmgreenbox}

\mdfdefinestyle{mdorangebox}{linecolor=RedOrange,backgroundcolor=RedOrange!5,
	linewidth=2pt,rightline=false,leftline=true,topline=false,bottomline=false,}
\declaretheoremstyle[headfont=\bfseries\sffamily\color{RedOrange!70!black},
mdframed={style=mdorangebox},headpunct={ --- },]{thmorangebox}

\mdfdefinestyle{mdblackbox}{linecolor=black,backgroundcolor=RedViolet!5!gray!5,
	linewidth=3pt,rightline=false,leftline=true,topline=false,bottomline=false,}
\declaretheoremstyle[mdframed={style=mdblackbox}]{thmblackbox}

\declaretheoremstyle[spaceabove=6pt,spacebelow=6pt]{basehead}

\declaretheorem[style=basehead,name=Problem]{problem}
\declaretheorem[style=basehead,name=Problem,numbered=no]{problem*}
\declaretheorem[style=thmbloobox,name=Setup]{setup} % for config geo
\declaretheorem[style=thmredbox,name=Example,sibling=problem]{example}
\declaretheorem[style=thmredbox,name=$\bigstar$ Problem,sibling=problem]{niceprob}
\declaretheorem[style=thmbluebox,name=Theorem,sibling=problem]{theorem}
\declaretheorem[style=thmbluebox,name=Corollary,sibling=theorem]{corollary}
\declaretheorem[style=thmbluebox,name=Proposition,sibling=theorem]{prop}
\declaretheorem[style=thmbluebox,name=Lemma,numberwithin=problem]{lemma}
\declaretheorem[style=thmbluebox,name=Theorem,numbered=no]{theorem*}
\declaretheorem[style=thmbluebox,name=Lemma,numbered=no]{lemma*}
\declaretheorem[style=thmgreenbox,name=Claim,numberwithin=problem]{claim}
\declaretheorem[style=thmgreenbox,name=Claim,numbered=no]{claim*}
\declaretheorem[style=thmblackbox,name=Remark,numberwithin=problem]{remark}
\declaretheorem[style=thmblackbox,name=Remark,numbered=no]{remark*}
\declaretheorem[style=thmgreenbox,name=Definition,sibling=theorem]{definition}
\declaretheorem[style=thmgreenbox,name=Definition,numbered=no]{definition*}
\declaretheorem[style=thmorangebox,name=Warning,sibling=theorem]{warning}
\declaretheorem[style=thmorangebox,name=Warning,numbered=no]{warning*}
\declaretheorem[style=thmorangebox,name=Note,numbered=no]{note}
\declaretheorem[style=thmblackbox,name=Exercise,sibling=problem]{exercise}
\declaretheorem[style=thmblackbox,name=Exercise,numbered=no]{exercise*}
\declaretheorem[style=thmblackbox,name=Question,sibling=problem]{question}
\declaretheorem[style=thmblackbox,name=Question,numbered=no]{question*}

\addtolength{\textheight}{3.14cm}
\providecommand{\ol}{\overline}
\providecommand{\ul}{\underline}
\providecommand{\eps}{\varepsilon}
\providecommand{\half}{\frac{1}{2}}
\providecommand{\dang}{\measuredangle} %% Directed angle
\providecommand{\CC}{\mathbb C}
\providecommand{\FF}{\mathbb F}
\providecommand{\NN}{\mathbb N}
\providecommand{\QQ}{\mathbb Q}
\providecommand{\RR}{\mathbb R}
\providecommand{\ZZ}{\mathbb Z}
\providecommand{\dg}{^\circ}
\providecommand{\ii}{\item}
\providecommand{\setdiff}{\smallsetminus}
\providecommand{\alert}[1]{{\sffamily\textbf{\textcolor{blue}{#1}}}}
\providecommand{\inv}{^{-1}}
\providecommand{\mb}[1]{\mathbf{#1}}
\providecommand{\dif}{\;d}
\newcommand*{\horzbar}{\rule[.5ex]{2.5ex}{0.5pt}}

\reversemarginpar

\newenvironment{soln}{\begin{proof}[Solution]}{\end{proof}}
\newenvironment{walkthrough}{\noindent \textbf{Walkthrough.}}{\medskip}
\newlist{walk}{enumerate}{3}
\setlist[walk]{label=\bfseries (\alph*)}

\providecommand{\pow}{\operatorname{Pow}}
\providecommand{\vocab}[1]{{\textbf{\textcolor{ForestGreen}{#1}}}}
\providecommand{\lcm}{\operatorname{lcm}}
\providecommand{\abs}[1]{\left|#1\right|}

\usepackage{asymptote}
\begin{asydef}
	defaultpen(fontsize(10pt));
	unitsize(1cm);
	size(8cm); // set a reasonable default
	usepackage("amsmath");
	usepackage("amssymb");
	settings.tex="pdflatex";
	settings.outformat="pdf";
	// Replacement for olympiad+cse5 which is not standard
	import geometry;
	// recalibrate fill and filldraw for conics
	void filldraw(picture pic = currentpicture, conic g, pen fillpen=defaultpen, pen drawpen=defaultpen)
	{ filldraw(pic, (path) g, fillpen, drawpen); }
	void fill(picture pic = currentpicture, conic g, pen p=defaultpen)
	{ filldraw(pic, (path) g, p); }
	// some geometry
	pair foot(pair P, pair A, pair B) { return foot(triangle(A,B,P).VC); }
	pair orthocenter(pair A, pair B, pair C) { return orthocentercenter(A,B,C); }
	pair centroid(pair A, pair B, pair C) { return (A+B+C)/3; }
	// cse5 abbreviations
	path CP(pair P, pair A) { return circle(P, abs(A-P)); }
	path CR(pair P, real r) { return circle(P, r); }
	pair IP(path p, path q) { return intersectionpoints(p,q)[0]; }
	pair OP(path p, path q) { return intersectionpoints(p,q)[1]; }
	path Line(pair A, pair B, real a=0.6, real b=a) { return (a*(A-B)+A)--(b*(B-A)+B); }
	// cse5 more useful functions
	picture CC() {
		picture p=rotate(0)*currentpicture;
		currentpicture.erase();
		return p;
	}
	pair MP(Label s, pair A, pair B = plain.S, pen p = defaultpen) {
		Label L = s;
		L.s = "$"+s.s+"$";
		label(L, A, B, p);
		return A;
	}
	pair Drawing(Label s = "", pair A, pair B = plain.S, pen p = defaultpen) {
		dot(MP(s, A, B, p), p);
		return A;
	}
	path Drawing(path g, pen p = defaultpen, arrowbar ar = None) {
		draw(g, p, ar);
		return g;
	}
\end{asydef}

\usepackage{mathrsfs}

\ihead{\footnotesize MAT 362 HW03}
\ohead{\footnotesize Last compiled \today}

\title{\vspace{-2em}MAT 362 HW03}
\author{\large Aditya Pahuja}
\date{\large Due 02/20}
\begin{document}
\maketitle

\begin{problem*}2.2.1.
	Show that the cylinder $\{(x,y,z)\in\RR^3 : x^2 + y^2 = 1\}$
	is a regular surface, and find parametrizations
	whose coordinate neighborhoods cover it.
\end{problem*}

\begin{soln}
	Consider the four maps $\mb x_1, \mb x_2,\mb x_3,\mb x_4\colon
	(-1,1)\times\RR\to\RR^3$ defined as follows:
	\begin{align*}
		\mb x_1(x,z) &= (x,\sqrt{1-x^2},z)\\
		\mb x_2(x,z) &= (x,-\sqrt{1-x^2},z)\\
		\mb x_1(y,z) &= (\sqrt{1-y^2},y,z)\\
		\mb x_2(y,z) &= (\sqrt{1-x^2},y,z).
	\end{align*}
	
	The union of the images of the $\mb x_i$ is the cylinder itself,
	the images of $\mb x_1$ and $\mb x_2$ together cover all points except
	the ones with a $y$-coordinate of zero, while $\mb x_3$ and $\mb x_4$
	cover the points on the cylinder with nonzero $x$-coordinate.
	
	We will verify that $\mb x_1$ is a parametrization of
	its image $\mb x_1(\RR^2)$; the other verifications are analogous,
	so this will prove that the cylinder is a regular surface.
	
	First, it's clear that all the component functions of $\mb x_1$
	are smooth; in particular $\sqrt{1-x^2}$ is a composition of
	smooth functions so it is smooth.
	
	Second, $\mb x_1$ is a homeomorphism between $\RR^2$
	and its image $\mb x_1(\RR^2)$, since it is clearly continuous
	and it has a continuous inverse
	in the form of a projection onto the $xz$-plane.
	
	Third, its differential map is the matrix
	\[\begin{pmatrix}
		1&0\\
		*&0\\
		0&1
	\end{pmatrix}\]
	where $*$ is some function of $x$ that we don't care about;
	this matrix has rank two because the rows $(1, 0)$ and $(0, 1)$
	are unconditionally independent of one another.
\end{soln}

\begin{problem*}2.2.3.
	Show that the two-sheeted cone with its vertex at the origin
	(that is, the set $S = \{(x,y,z)\in\RR^3 : x^2+y^2-z^2=0\}$)
	is not a regular surface.
\end{problem*}

\begin{soln}
	Suppose that $\mb x\colon U\to V\cap S$ is a parametrization
	that defines a coordinate neighborhood
	$V\cap S$ of $0$. Then, observe that $(V\cap S)\setdiff\{0\}$
	is disconnected because $S\setdiff\{0\}$ is disconnected.
	We show that this implies $S$ cannot be regular.
	
	\begin{claim*}
		If $U$ is an open set in $\RR^3$,
		then it has no cut points;
		that is, there are no points $x\in U$ such that
		$U\setdiff\{x\}$ is disconnected.
	\end{claim*}
	\begin{proof}
		Suppose $U$ has a cut point $x$ and let $U' = U\setdiff\{x\}$.
		Then, there exist nonempty open sets $A,B\subseteq\RR^3$
		that do not intersect in $U'$ and whose union contains $U'$.
		Let $B_\eps(x)\subseteq U$ be an open ball centered at $x$
		with radius $\eps$ contained inside $U$. It is easy to show that
		$B_\eps(x)\setdiff\{x\}$ is connected; in fact, it is path connected,
		for the ball (with $x$) is convex, and so two points in
		the ball without $x$ can be connected by a segment, or,
		if the segment passes through $x$, then by a path made by
		joining two segments, with the intermediary endpoint
		not being collinear with the endpoints of the path.
		
		Since $B_\eps(x)\setdiff\{x\}$ is connected,
		it must lie entirely in one of $A$ or $B$ --- say WLOG $A$ --- else
		$A$ and $B$ form a separation of $B_\eps(x)\setdiff\{x\}$,
		contradicting its connectedness.
		But then the set $A' = A\cup \{x\}$ is open
		because it is the union of $A$ with $B_\eps(x)$
		and doesn't intersect $B$ inside $U$, so $A$ and $B$
		exhibit that $U$ is disconnected, which is a contradiction.
	\end{proof}

	Since $U$ and $V\cap S$ are homeomorphic, a cut point of $U$
	maps to a cut point of $V\cap S$ and vice versa ---
	but $U$ has no cut points, so $V\cap S$ must have no cut points.
	This is false in our situation here, so $S$ cannot be regular.
\end{soln}

\begin{problem*}2.2.7.
	Let $f(x,y,z) = (x + y + z - 1)^2$.
	\begin{enumerate}[(a)]
		\ii Locate the critical points and critical values of $f$.
		\ii For what values of $c$ is the set $f(x,y,z) = c$
		a regular surface?
	\end{enumerate}
\end{problem*}

\begin{soln}
	The differential of $f$ is
	\[\mathrm df_{(x,y,z)} = 2(x+y+z-1)\cdot(1,1,1).\]
	Thus $(x,y,z)$ is a critical point iff $x+y+z-1 = 0$.
	The corresponding critical value is just $f(x,y,z) = 0$.
	
	Since $f\inv(c)$ is a regular surface
	whenever $c$ is a regular value, any $c\ne 0$ makes the set
	a regular surface.
	
	We then have to check whether $f(x,y,z) = 0$ is a regular surface.
	This is trivial; the set $f\inv(0)$ is $\{(x,y,z)\in\RR^3:x+y+z=1\}$
	which is a plane and thus regular,
	for example with the parametrization $(u,v)\mapsto(u,v,1-u-v)$.
	Clearly this is infinitely differentiable and a homeomorphism
	between $\RR^2$ and the plane, and the derivative
	\[\begin{pmatrix}
		1&0\\
		0&1\\
		-1&-1
	\end{pmatrix}\]
	is obviously injective.
\end{soln}

\begin{problem*}2.2.16.
	One way to define a system of coordinates for the sphere $S^2$,
	given by $x^2 + y^2 + (z-1)^2 = 1$, is to consider
	the so-called \vocab{stereographic projection}
	$\pi\colon S\setdiff\{N\}\to\RR^2$ which carries a point
	$p\in S^2$ that is not equal to the north pole $N = (0,0,2)$
	to the intersection of the $xy$-plane with the line through $N$ and $p$.
	Let $(u,v) = \pi(x,y,z)$ where $(x,y,z)\in S^2\setdiff\{N\}$.
	\begin{enumerate}[(a)]
		\ii Show that $\pi\inv\colon\RR^2\to S^2$ is given by
		\[\pi\inv(u, v) = \begin{pmatrix}
			\dfrac{4u}{u^2 + v^2 + 4}\\[2ex]
			\dfrac{4v}{u^2 + v^2 + 4}\\[2ex]
			\dfrac{2(u^2 + v^2)}{u^2 + v^2 + 4}
		\end{pmatrix}.\]
		\ii Show that it is possible, using stereographic projection,
		to cover the sphere with two coordinate neighborhoods.
	\end{enumerate}
\end{problem*}

\begin{soln}
	(a) Let $O = (0,0,0)$ be the origin, $A = (u,v,0)$ be
	the point in the $xy$-plane, and $B = \pi\inv(u,v)$
	be the point we are trying to find, namely the second intersection of
	line $\ol{NA}$ with the sphere. Then, since $\ol{ON}$
	is a diameter of the sphere, $\angle NBO = \frac\pi2$,
	so we are trying to compute the foot of the altitude from
	$O$ to line $\ol{NA}$. The vector $\overrightarrow{OB}$ can then
	simply be computed as
	\begin{align*}
		\overrightarrow{ON} + \overrightarrow{NB}
		&= (0,0,2) +
		\frac{\overrightarrow{NO}\cdot\overrightarrow{NA}}
		{\|\overrightarrow{NA}\|}
		\cdot\frac{\overrightarrow{NA}}{\|\overrightarrow{NA}\|}\\
		&= (0,0,2) + \frac{(0,0,-2)\cdot(u,v,-2)}{\|(u,v,-2)\|^2}
		\cdot (u,v,-2)\\
		&= (0,0,2) + \frac{4}{u^2+v^2+4}\cdot (u,v,-2)\\
		&= \frac{1}{u^2+v^2+4}(4u,4v,-8+2(u^2+v^2+4))\\
		&= \frac{1}{u^2+v^2+4}(4u,4v,2(u^2+v^2)).
	\end{align*}
	
	(b) Let $\pi_1 \equiv \pi$ be the stereographic projection
	of $S^2\setdiff{N}$ onto the $xy$-plane and let
	$\pi_2$ be the stereographic projection of $S^2\setdiff\{O\}$
	onto the plane $z = 2$, where $O = (0,0,0)$ is the south pole.
	To be more precise $\pi_2$ projections the sphere onto $z=2$
	and then forgets the $z$-coordinate/projects $z=2$ to the $xy$-plane.
	I claim that these two maps exhibit a covering of $S^2$
	with the two coordinate neighborhoods $S^2\setdiff\{N\}$
	and $S^2\setdiff\{O\}$.
	
	We can see clearly that $\pi_1\inv\colon\RR^2\to S^2\setdiff\{N\}$
	and $\pi_2\colon\RR^2\to S^2\setdiff\{O\}$ are
	differentiable of all orders (the only possible obstruction
	would be creating a denominator of zero somehow, but the denominators
	of the rational functions would just be powers of $u^2+v^2+4$,
	which are always nonnegative), so this verifies condition (1).
	
	Note now that $\pi_1\inv$ and $\pi_2\inv$ are
	bijections (since they are invertible, having inverses $\pi_1$, $\pi_2$,
	by part (a)) from an open subset of $\RR^2$
	to an open subset $S^2$ (where ``open'' is with respect to
	the subspace $S^2$, i.e. the intersection of an open set
	in $\RR^2$ with $S^2$). To check condition (2) we need to also check
	that $\pi_1$, $\pi_2$ are continuous. We will just check it for $\pi_1$
	by directly computing $\pi_1(x,y,z)$. If we are given
	\[\pi\inv(u,v) = (x,y,z),\]
	then we can compute
	\[z = \frac{2(u^2+v^2)}{u^2+v^2 + 4}\iff
	u^2+v^2 = \frac{4z}{2-z}\]
	(noting that $2-z\ne 0$ because $z < 2$) and then
	\[x = \frac{4u}{4z/(2-z) + 4} = \frac{4u}{8/(2-z)} = \frac{u(2-z)}{2}
	\iff u = \frac{2x}{2 - z}\]
	and similarly $v = \frac{2y}{2-z}$. These are clearly continuous
	in $x,y,z$, so $\pi_1$ is continuous. A similar proof follows for $\pi_2$
	(or one can apply the change of coordinates $z\mapsto 2-z$
	to reflect the stereographic projection over $z=1$,
	swapping $O$ and $N$ as well as the planes that are being mapped to).
	We are hence done verifying condition (2).
	
	To finish proving that these are coordinate neighborhoods we need to check
	that the derivatives of $\pi_1\inv$ and $\pi_2\inv$ are injective.
	Again we check this only in the case of $\pi_1\inv$: first we compute
	\[\mathrm d\pi_1\inv_{(u,v)} =
	\frac1{(u^2+v^2+4)^2}
	\begin{pmatrix}
		4(-u^2+v^2+4) & -8uv\\
		-8uv & 4(u^2-v^2+4)\\
		16u & 16v
	\end{pmatrix}.\]
	To check that this linear map is one-to-one,
	we need the three equations
	\begin{align*}
		4(-u^2 + v^2 + 4) &= -8\lambda uv\\
		-8uv &= 4\lambda(u^2 - v^2 + 4)\\
		16u &= 16\lambda v
	\end{align*}
	to simultaneously hold for some real $\lambda \ne 0$,
	noting that neither of the two columns can be all zero at the same time,
	since either $u$ or $v$ would have to be zero
	but then $4(-u^2+v^2+4)$ or $4(u^2-v^2+4)$ respectively would be nonzero.
	This also means we can assume $u,v\ne0$.
	We can now get $\lambda = \frac uv$ so that
	\begin{align*}
		4(-u^2 + v^2 + 4) &= -8u^2\\
		4(u^2 - v^2 + 4) &= -8v^2.
	\end{align*}
	In fact both equations are equivalent to $4u^2 + 4v^2 = -16$,
	which is impossible because squares are nonnegative.
	This checks off condition (3) for $\pi_1\inv$,
	and by an analogous argument or by symmetry, we can conclude
	that it holds for $\pi_2\inv$ as well, thus finishing the problem.
\end{soln}
\pagebreak
\begin{problem*}2.3.1.
	Let $S^2$ be the unit sphere and let
	$A\colon S^2\to S^2$ be the \vocab{antipodal map}
	$A(x,y,z) = -(x,y,z)$. Prove that $A$ is a diffeomorphism.
\end{problem*}

\begin{soln}
	Since $A$ is its own inverse, we only need to show that
	$A$ is differentiable on all of $S^2$.
	By definition, we wish to show that for all $x\in A$
	and for any two parametrizations
	$\mb x_1\colon U_1\to S^2$, $\mb x_2\colon U_2\to S^2$
	such that $x\in\mb x_1(U_1)\cap\mb x_2(U_2)$,
	the composition
	\[\mb x_2\inv\circ A\circ\mb x_1\]
	is differentiable at $\mb x_1\inv(x)$ in the sense that
	a map from $\RR^2$ to $\RR^2$ is differentiable. This is immediate
	in light of the fact that $-v\mapsto v$ is a differentiable map
	on all of $\RR^3$, and compositions of differentiable functions
	are differentiable.
\end{soln}

\begin{problem*}2.3.2.
	Let $S\subseteq\RR^3$ be a regular surface and let
	$\pi\colon S\to\RR^2$ be the map that takes each $p\in S$
	into its orthogonal projection over $\RR^2 = \{(x,y,z)\in\RR^3 : z=0\}$.
	Is $\pi$ differentiable?
\end{problem*}

\begin{soln}
	We observe that $\pi(x,y,z) = (x,y,0)$ is differentiable
	when treated as a function from $\RR^3$ to $\RR^2$.
	Thus, if $\mb x\colon U\to V$ is a parametrization of
	a coordinate neighborhood of some $p\in S$,
	then, noting that $\mb x$ is differentiable on $U$
	and $\pi$ is differentiable on $\mb x(U) = V$
	(in the $\RR^3\to\RR^2$ sense),
	$\pi\circ \mb x\colon U\to\RR^2$ is differentiable on all of $U$,
	hence $\pi$ is differentiable.
\end{soln}

\begin{problem*}2.3.4.
	Construct a diffeomorphism between the ellipsoid
	\[\frac{x^2}{a^2} + \frac{y^2}{b^2} + \frac{z^2}{c^2} = 1\]
	and the sphere $x^2 + y^2 + z^2 = 1$.
\end{problem*}

\begin{soln}
	Let $E$ be the ellipsoid and $S^2$ the sphere.
	Then, consider the map $f\colon E\to S^2$ given by
	$f(x,y,z) = (x/a,y/b,z/c)$. Clearly this is a bijection
	between points on $E$ and $S^2$, for $f(x,y,z)$ is in $S^2$ iff
	\[(x/a)^2 + (y/b)^2 + (z/c)^2 = 1\]
	which holds iff $(x,y,z)$ is on $E$, and this map has an inverse
	$f\inv(x,y,z) = (ax,by,cz)$.
	Moreover, $f$ is differentiable because it is the restriction of
	a differentiable function $\RR^3\to\RR^3$ to $E$
	and its inverse is differentiable because it is similarly
	the restriction of a differentiable function $\RR^3\to\RR^3$ to $S^2$.
	Thus $f$ is a diffeomorphism between $E$ and $S^2$.
\end{soln}

\begin{problem*}2.3.5.
	Let $S\subseteq\RR^3$ be a regular surface, and let
	$d\colon S\to\RR$ be given by $d(p) = |p - p_0|$, where
	$p\in S$, $p_0\in\RR^3$, $p_0\notin S$.
	Prove that $d$ is differentiable.
\end{problem*}

\begin{soln}
	We show that $d'\colon\RR^3\to\RR$ defined by
	$d'(p) = |p - p_0|$ (so $d'|_S = d$) is differentiable
	at all points besides $p_0$, at which point the result follows
	because the restriction of a map $\RR^3\to\RR$
	that is differentiable at $x$ to a map $S\to\RR$
	is still differentiable at $x$.
	
	We first show that $d'(p)^2$ is differentiable on all of $\RR^3$;
	this follows from $d'(p)^2 = |p - p_0|^2 = (p - p_0)\cdot(p - p_0)$
	and so its derivative is the linear map given by the $1\times 3$ matrix
	$2(p - p_0)$ (so the derivative at a point $x$ is $2(p-p_0)\cdot x$).
	Then, $\sqrt{d'(p)^2}$ is differentiable everywhere except at $p_0$
	because $\sqrt\bullet\colon\RR_{\ge0}\to\RR$
	is differentiable on all positive reals but not at zero.
	
	By the first paragraph, this finishes the problem.
\end{soln}























\end{document}